
% ============================================================
%  INTRODUCTION GÉNÉRALE
% (pagenumbering/setcounter are in main.tex)
% ============================================================
\cleardoublepage

\chapter*{Introduction Générale}
\addcontentsline{toc}{chapter}{Introduction Générale}
\markboth{Introduction Générale}{Introduction Générale}

L'agriculture mondiale connaît aujourd'hui une transformation profonde sous
l'impulsion de la \textit{quatrième révolution industrielle}. La convergence de
l'Internet des Objets (IoT), de l'intelligence artificielle (IA), des systèmes
d'information géographique (SIG) et des architectures logicielles distribuées a
favorisé l'émergence de l'\textbf{agriculture intelligente}, définie comme une agriculture de
précision, connectée et pilotée par la donnée. Cette évolution répond à des
enjeux majeurs de productivité, de durabilité et de résilience. Selon la
FAO~\cite{fao2023}, la population mondiale devrait atteindre 9,7 milliards
d'individus en 2050, ce qui impose une augmentation significative de la
production alimentaire mondiale. Dans ce contexte, l'intégration de l'IA dans
les systèmes agricoles apparaît comme un levier stratégique pour améliorer la
prise de décision, optimiser l'utilisation des ressources et renforcer la
capacité d'adaptation des exploitations face aux contraintes climatiques,
économiques et sanitaires.

Ces enjeux prennent une importance particulière dans le contexte tunisien. En
effet, le secteur agricole représente environ 10\% du PIB national et emploie
près de 16\% de la population active (UTAP~\cite{utap2023}). Malgré son rôle
économique et social essentiel, ce secteur demeure confronté à plusieurs défis
structurels, notamment le morcellement foncier, le vieillissement de la
population agricole, l'accès limité aux technologies numériques et la dépendance
vis-à-vis de solutions étrangères souvent coûteuses ou insuffisamment adaptées
aux réalités locales. À ces contraintes s'ajoute l'absence d'outils numériques
conçus en langue locale, en particulier en Darija tunisienne, ce qui limite
l'adoption technologique par les agriculteurs et les éleveurs. La conception de
solutions souveraines, accessibles, contextualisées et capables de fonctionner
dans des environnements à connectivité limitée constitue ainsi un enjeu central
pour la modernisation du secteur agricole tunisien.

C'est dans cette perspective que s'inscrit le projet \textbf{Smart Farm AI
Enterprise}. Il s'agit d'une plateforme d'agriculture intelligente conçue pour
répondre aux besoins opérationnels des exploitations agricoles et d'élevage. La
solution couvre la gestion multi-espèces, incluant les bovins, ovins, caprins,
volailles et abeilles, tout en intégrant des modules d'analyse d'image par
vision par ordinateur, un assistant IA conversationnel en Darija tunisienne, une
surveillance IoT en temps réel et un ERP avicole complet. L'ensemble repose sur
une architecture multilocataire sécurisée et sur une application mobile de type
PWA conçue selon une approche \textit{hors connexion d'abord}, afin de garantir l'accessibilité du système même
dans les zones rurales où la connexion Internet demeure instable.

Le présent mémoire a pour objectif de documenter la démarche de conception,
d'implémentation et de validation de cette plateforme. Il vise, d'une part, à
présenter l'architecture logicielle adoptée pour intégrer de manière cohérente
les composantes IA, IoT et web, et, d'autre part, à évaluer les performances des
modèles déployés à l'aide de méthodes statistiques rigoureuses, telles que les
intervalles de confiance et la validation croisée. Il s'attache également à
analyser la validité interne et externe des résultats obtenus, tout en proposant
une contribution à l'état de l'art sur les systèmes combinant RAG et LLM dans
une langue dialectale non standard, en l'occurrence la Darija tunisienne
appliquée au domaine de l'élevage.

La démarche de recherche repose sur trois hypothèses principales, vérifiées
expérimentalement dans les chapitres suivants. La première hypothèse concerne
l'efficacité d'une chaîne de traitement RAG+LLM locale, fondée sur Labess-7B, pour traiter des
requêtes agricoles formulées en Darija tunisienne ; les résultats obtenus
montrent un accord inter-annotateurs de $\kappa = 0.76$, supérieur au seuil de
$\kappa = 0.60$ généralement associé à un accord substantiel selon Landis et
Koch~\cite{landis1977}. La deuxième hypothèse porte sur la stabilité des modèles
YOLOv11 entraînés avec différentes graines aléatoires ; l'écart-type mesuré du
mAP@50 global est de $\sigma = 0.45\%$, confirmant une variabilité inférieure à
$\pm 1\%$. Enfin, la troisième hypothèse évalue l'adoption d'une application
mobile conçue selon une approche \textit{hors connexion d'abord} par les travailleurs agricoles formés ; le taux
d'adoption atteint 91\% après une semaine d'usage sur les sept fermes pilotes,
dépassant ainsi le seuil initialement fixé à 80\%.

Au-delà de son apport applicatif, ce travail présente plusieurs contributions
originales. Il propose d'abord, à la connaissance de l'auteur, l'une des premières
évaluations quantitatives d'un système RAG+LLM en Darija agricole, combinant des
mesures telles que BLEU, ROUGE et le coefficient Kappa. Il met ensuite en place
un protocole de validation statistique rarement appliqué à un ERP agricole
opérationnel, notamment à travers l'utilisation d'une validation croisée à cinq plis
avec \textit{TimeSeriesSplit}, d'intervalles de confiance à 95\% et d'une analyse
des menaces à la validité. Enfin, il propose une architecture agrotechnologique
à code source ouvert et souveraine, intégrant 11 catégories YOLO, 5 modèles ML avicoles,
RAG ChromaDB et un LLM local dans un déploiement Docker reproductible, adapté
aux pays émergents et aux environnements à connectivité limitée.

Le mémoire est structuré en quatre chapitres complémentaires. Le
\textbf{Chapitre~1} présente l'organisme d'accueil, le contexte général du projet
et la méthodologie de gestion hybride CRISP-DM/Scrum. Le \textbf{Chapitre~2}
détaille l'analyse des besoins ainsi que l'architecture globale du système. Le
\textbf{Chapitre~3} expose la modélisation des composants d'intelligence
artificielle et l'évaluation quantitative de leurs performances. Enfin, le
\textbf{Chapitre~4} décrit la réalisation technique de la plateforme et sa
validation industrielle sur le terrain.

% ============================================================
%  CHAPITRE 1 : CONTEXTE GÉNÉRAL ET CADRE MÉTHODOLOGIQUE
% ============================================================
